% Vietnamese translation for OI Public Library
% Source-File: IOI中国国家候选队论文/2008/Day2/8.高亦陶《从立体几何问题看降低编程复杂度》/从立体几何问题看降低编程复杂度.ppt
\documentclass[11pt]{article}
\usepackage{oipl-vn}
\usepackage{tikz}
\usetikzlibrary{calc}

\newcommand{\Oh}{\mathcal{O}}

\title{Bản trình chiếu: Giảm độ phức tạp lập trình qua bài hình học không gian}
\author{Gao Yitao}
\date{2008}

\begin{document}
\maketitle

\origtitle{Reducing implementation complexity through solid geometry}
\sourcepath{IOI China candidate team papers / 2008 / Day 2 / Gao Yitao.ppt}

\section{Phạm vi bản dịch}

Tệp gốc là PowerPoint-only. Dòng chữ trích được từ PPT khá ít, nhưng các mốc
ổn định gồm \(\Oh(1)\), \textit{Model Rocket Height}, \textit{Crossing Prisms},
khoảng \([0,10]\), mặt cắt \(y=y_k\), công thức có \(\cos\), và các tên hình
minh họa như \texttt{mrh\_angle}, \texttt{cpr\_s1}, \texttt{cpr\_w1}. Bản ghi
chú này dịch và hệ thống hóa mạch slide theo các mốc đó; những hình vẽ trang
trí hoặc hoạt ảnh được gộp thành mô tả toán học.

Chủ đề không phải là trình bày một thư viện hình học 3D tổng quát. Ngược lại,
slide nhấn mạnh cách chọn mô hình để bài toán không gian được rút về vài phép
tính một chiều hoặc hai chiều, nhờ đó cài đặt ngắn hơn và ít trường hợp biên
hơn.

\section{Động cơ}

Trong hình học không gian, độ khó lập trình thường đến từ các chi tiết:
vector ba chiều, giao đường thẳng với mặt phẳng, thứ tự điểm trên đa giác cắt,
sai số số thực và các cấu hình suy biến. Nếu mô hình quá tổng quát, lời giải
dễ biến thành một tập lớn các hàm hình học khó kiểm soát.

Vì vậy slide đặt mục tiêu ``giảm độ phức tạp lập trình'': tìm phép chiếu, mặt
cắt, hệ tọa độ hoặc công thức lượng giác để phần 3D biến thành bài toán đã
quen. Khi đại lượng cần tính chỉ phụ thuộc vào một vài góc, độ dài hoặc lát
cắt, ta có thể tránh hầu hết thao tác hình học không gian trực tiếp.

\section{\textit{Model Rocket Height}: đổi bài 3D thành lượng giác}

Ví dụ đầu là \textit{Model Rocket Height}. Các slide có nhiều hình
\texttt{mrh1b}, \texttt{mrh2b}, \texttt{mrh\_angle} và \texttt{mrh\_f}, cho
thấy mạch lập luận xoay quanh việc dựng tam giác quan sát và đọc chiều cao
từ các góc.

\begin{center}
\begin{tikzpicture}[scale=0.9, >=stealth]
  \draw[thick] (-0.2,0) -- (6.2,0);
  \coordinate (A) at (0.5,0);
  \coordinate (B) at (4.8,0);
  \coordinate (R) at (3.3,3.1);
  \draw[fill=black] (A) circle (1.5pt) node[below] {trạm 1};
  \draw[fill=black] (B) circle (1.5pt) node[below] {trạm 2};
  \draw[fill=black] (R) circle (1.5pt) node[above] {tên lửa};
  \draw[dashed] (R) -- ($(R)!(A)!(B)$) node[midway,right] {$h$};
  \draw[thick] (A) -- (R) -- (B);
  \draw (1.2,0) arc[start angle=0,end angle=37,radius=0.7];
  \node at (1.45,0.27) {$\alpha$};
  \draw (4.15,0) arc[start angle=180,end angle=132,radius=0.65];
  \node at (3.9,0.35) {$\beta$};
\end{tikzpicture}
\end{center}

Thay vì dựng đường thẳng 3D và tìm giao chính xác, ta chọn mặt phẳng chứa hai
trạm quan sát và tên lửa. Trong mặt phẳng này, bài toán chỉ còn là tam giác:
biết khoảng cách giữa hai trạm và các góc nâng, tính chiều cao \(h\). Nếu gọi
khoảng cách từ chân đường cao tới trạm thứ nhất là \(x\), thì
\[
  h=x\tan\alpha=(d-x)\tan\beta.
\]
Suy ra
\[
  h=\frac{d\tan\alpha\tan\beta}{\tan\alpha+\tan\beta}.
\]
Trong các biến thể có thêm góc phương vị, bước quan trọng vẫn là chiếu về mặt
phẳng thích hợp trước, rồi mới dùng lượng giác phẳng. Đó là ý nghĩa của mốc
\(\Oh(1)\) trong slide: sau khi chọn đúng biểu diễn, mỗi trường hợp chỉ cần
một số hằng phép tính.

\section{\textit{Crossing Prisms}: thể tích bằng tích phân lát cắt}

Ví dụ thứ hai là \textit{Crossing Prisms}. Dòng chữ trích được nhắc khoảng
\([0,10]\), mặt cắt \(y=y_k\), các hình \texttt{cpr\_w1}, \texttt{cpr\_w2},
\texttt{cpr\_w3c}, \texttt{cpr\_w3d}, và công thức dạng
\[
  2\sum \frac{\cdots}{\cos(\cdots)}.
\]
Mạch slide vì vậy là phương pháp lát cắt: thay vì xử lý trực tiếp giao của
hai lăng trụ xiên, ta cắt vật thể bởi các mặt phẳng \(y=y_k\), tính diện tích
mặt cắt, rồi cộng hoặc tích phân các lát.

\begin{center}
\begin{tikzpicture}[scale=0.85, >=stealth]
  \draw[->] (-0.4,0) -- (6.4,0) node[right] {$x$};
  \draw[->] (0,-0.3) -- (0,3.7) node[above] {$z$};
  \draw[thick,blue] (0.4,0.4) -- (5.8,0.4) -- (4.6,3.0) -- (1.0,3.0) -- cycle;
  \draw[thick,red] (0.9,0.1) -- (5.1,3.2);
  \draw[thick,red] (1.8,0.1) -- (6.0,3.2);
  \draw[dashed] (0.2,1.6) -- (6.1,1.6);
  \node[left] at (0.2,1.6) {$y=y_k$};
  \node[blue] at (5.2,2.8) {lăng trụ 1};
  \node[red] at (5.4,1.5) {lăng trụ 2};
\end{tikzpicture}
\end{center}

Ở mỗi lát \(y=y_k\), giao của hai lăng trụ trở thành giao của vài đoạn hoặc
đa giác trong mặt phẳng. Các đoạn biên có thể dài hơn hình chiếu của chúng
một hệ số \(1/\cos\theta\); vì vậy công thức tổng trong slide xuất hiện mẫu
\(\cos\). Khi các điểm gãy theo \(y\) được sắp thứ tự, diện tích mặt cắt là
một hàm tuyến tính hoặc bậc thấp trên từng khoảng, nên thể tích có thể tính
bằng cộng các phần đơn giản.

Điểm cần học không phải là một công thức duy nhất cho mọi lăng trụ, mà là cách
đổi bài toán: từ giao khối 3D khó cài đặt sang một chuỗi bài toán mặt cắt 2D
có thứ tự rõ ràng. Các ký hiệu như \(1,2+,2-,3\) và \(3,2-,2+,1\) trong text
stream nhiều khả năng là thứ tự các biên khi quét qua các điểm gãy; đây chính
là phần giúp tránh xét tay quá nhiều trường hợp.

\section{Nguyên tắc cài đặt}

Từ hai ví dụ, slide rút ra một số nguyên tắc thực dụng:
\begin{itemize}
  \item Luôn hỏi đại lượng cần tính có thật sự cần mô hình 3D đầy đủ hay không.
  \item Nếu đề bài chỉ hỏi chiều cao, góc, thể tích hoặc diện tích cắt, hãy
  ưu tiên chiếu hoặc lát cắt.
  \item Đưa các cạnh xiên về hình chiếu và hệ số lượng giác như \(1/\cos\theta\)
  thay vì thao tác trực tiếp với mặt phẳng tổng quát.
  \item Chọn tham số quét để thứ tự sự kiện ít thay đổi nhất; khi thứ tự đổi,
  tách thành một khoảng mới.
  \item Giảm số nhánh suy biến bằng cách làm việc với công thức trên đoạn/mặt
  cắt đã chuẩn hóa.
\end{itemize}

\section{Kết luận}

Các bài hình học không gian trong thi lập trình không nhất thiết phải được
giải bằng một bộ công cụ 3D lớn. Nếu chọn đúng phép quy đổi, \textit{Model
Rocket Height} trở thành một bài lượng giác phẳng \(\Oh(1)\), còn
\textit{Crossing Prisms} trở thành bài quét các mặt cắt \(y=y_k\). Giảm độ
phức tạp lập trình ở đây là giảm số chiều của suy nghĩ trước khi viết mã.

\end{document}
